\section{Objective}\label{sec:objective}

Now that the main developments that led to the rise of \gls{dl} have been outlined, it is necessary to formalize what is actually being achieved when performing \gls{dl}.
To ensure clarity, the supervised learning setting will be considered, specifically applied to a regression task.
This framework serves as a standard reference in the field, as all building blocks of \gls{dl} applications are built upon this fundamental formalism.

In this setting, learning consists of progressively shaping a \gls{nn} that serves as an approximation of a target function.
At initialization, the network does not yet provide a meaningful representation of this function.
Since a \gls{nn} is a parametric function, its parameters are progressively updated during training so as to make it a better approximation of the target function.
It gradually becomes capable of capturing the relationship between inputs and outputs through successive parameter updates during training.
For this reason, \gls{dl} is naturally situated within the framework of machine learning.
A standard definition of machine learning is:
\begin{quote}
\emph{``A computer program is said to learn from experience $E$ with respect to some class of tasks $T$ and performance measure $P$, if its performance at tasks in $T$, as measured by $P$, improves with experience $E$.''} \cite{mitchell1997machine}
\end{quote}

When this formulation is adapted to \gls{dl}, the following correspondences hold:
\begin{itemize}
\item The computer program corresponds to the learning algorithm, which is responsible for adjusting the network's parameters.
\item The task $T$ consists of learning a \gls{nn} that provides a good approximation of the target function.
\item The performance measure $P$ corresponds to a function used to evaluate the quality of this approximation.
\item The experience $E$ takes the form of a dataset composed of input--output pairs derived from the function to be approximated.
\end{itemize}

% The mathematical formalism used in this chapter is now introduced in order to remove any ambiguity in the developments that follow.
The function to be approximated is denoted by $f$.
As with any function, it defines a mapping from an input to an output: $x \in \mathcal{X}$ denotes an input and $y \in \mathcal{Y}$ the corresponding output.
The function $f$ is therefore characterized by an input space $\mathcal{X}$ and an output space $\mathcal{Y}$, and can be formally written as $f : \mathcal{X} \to \mathcal{Y}$.
For simplicity of exposition, both $\mathcal{X}$ and $\mathcal{Y}$ are assumed to be vector spaces throughout this manuscript.
This assumption is adopted solely to facilitate the presentation, as the framework can be readily extended to more general settings.
The notation $\hat{\cdot}$ indicates that a given object is an approximation of another.
Since the goal of the trained \gls{nn} is to approximate $f$, this approximation is denoted by $\hat{f} : \mathcal{X} \to \mathcal{Y}$.
Given an input $x \in \mathcal{X}$, the output of the approximating function $\hat{f}(x)$ is denoted by $\hat{y}$.

It is important to emphasize that the target function $f$ is generally not known explicitly.
If it were, there would be no need to learn an approximation, since the true function would already provide the exact mapping.
In practice, only a finite set of observations generated by $f$ is available, organized in what is called a dataset.
To distinguish between the different samples, an index $i$ is introduced, allowing input--output pairs of the form $(x^{(i)}, y^{(i)})$ to be defined.
This leads to the following formulation for a dataset $\mathcal{D}$ of size $N$, where $N$ denotes the number of available examples:
\begin{equation}
    \mathcal{D} = \{(x^{(i)}, y^{(i)})\}_{i=1}^{N}.
    \label{eq:dataset}
\end{equation}

Now that the representation of the target function $f$ has been established, the next question is what makes a function $\hat{f}$ a good approximation of it.
For a given input $x^{(i)}$, this means that the \gls{nn} output $\hat{y}^{(i)}$ should be close to the true output $y^{(i)}$.
To quantify the quality of this approximation, a function $\mathcal{L}$ is introduced that measures the discrepancy between predictions and target values.
This function called the loss function plays a central role as it guides the learning process.
In regression settings, a common choice is the \gls{mse}, defined as:
\begin{equation}
    \mathcal{L}(\hat{y}, y) = (\hat{y} - y)^2.
    \label{eq:mse}
\end{equation}
Thanks to this loss, it is now possible to rigorously define the performance measure $P$.
This performance measure corresponds to the average loss over the dataset and can therefore be written as:
\begin{equation}
    P = \frac{1}{N} \sum_{i=1}^{N} \mathcal{L}(\hat{f}(x^{(i)}), y^{(i)}).
    \label{eq:performance_measure}
\end{equation}

This formulation captures the mathematical core of the learning objective, namely finding a \gls{nn} $\hat{f}$ that minimizes the prediction error over the dataset $\mathcal{D}$.
Since the target function $f$ may take many different forms, $\hat{f}$ must be expressive enough to approximate a broad range of such functions.
\gls{nn} architectures are well suited to this requirement, as they can represent a wide range of functional mappings, whose composition is detailed in the following section.